% GENERATED FILE. Edit canonical lessons, then run npm run generate:books.
% canonical-source-hash: fnv1a64:75784213ec200b00
% canonical-lessons: TA-C487-koluppu, TA-C487-puratam, TA-C487-vaittamin, TA-C487-ceri, TA-C487-eppam-vitu

\chapter{Food and Digestion}
\label{ch:ta-food-and-digestion}
\hlchaptermodality{487}

\begin{chapteropening}
\textbf{By the end of this chapter:} \emph{I can talk about fat, protein, vitamins, digestion and burping.}

Complete the last lesson of chapter 487: I can talk about fat, protein, vitamins, digestion and burping.
\end{chapteropening}

\section[\texorpdfstring{\ta{கொழுப்பு} (koḻuppu)}{கொழுப்பு (koḻuppu)}]{\ta{கொழுப்பு} (koḻuppu) — body fat; cholesterol}

\label{lesson:TA-C487-koluppu}



\begin{quote}
Before the new one: say the Tamil for a walking stick, then the Tamil for to pull through.
\end{quote}

\subsection*{You'll want to know: \ta{கொழுப்பு}}
\textbf{\ta{கொழுப்பு}} — \emph{koḻuppu} — \textquotedblleft{}body fat; cholesterol\textquotedblright{}.

\begin{quote}
\ta{எண்ணெயில்} \ta{பொரித்த} \ta{உணவில்} \ta{கொழுப்பு} \ta{அதிகம்}.
\end{quote}

\begin{quote}
\emph{eṇṇeyil poritta uṇavil koḻuppu atikam.}
\end{quote}

\begin{quote}
Deep-fried food is high in fat.
\end{quote}

In speech also cheek: \ta{அவனுக்குக்} \ta{கொழுப்பு}, he is cocky.

\begin{cousinweb}
\textbf{\ta{கொழுப்பு}} (\emph{koḻuppu}) is native Dravidian, from \ta{கொழு} \emph{koḻu}, \textquotedblleft{}fat; to flourish\textquotedblright{} (DEDR 2146).
\end{cousinweb}

\begin{grammarlens}[title={What you've built}]
One more word for describing what you lived and felt.
\end{grammarlens}

\subsection*{Your turn}
\begin{itemize}
\raggedright
  \item \emph{Say it:} \emph{koḻuppu}
  \item \emph{Say it:} \emph{koḻuppu}, once more
  \item \emph{Say it:} \emph{eṇṇeyil poritta uṇavil koḻuppu atikam.} — Deep-fried food is high in fat.
  \item \emph{Recall:} say the Tamil for a walking stick, then the Tamil for to pull through, then say \emph{koḻuppu} again
\end{itemize}

\begin{tcolorbox}[breakable,colback=teal!4,colframe=teal!35!black,title={Before you move on}]
What does \ta{கொழுப்பு} mean? (\textquotedblleft{}Body fat; cholesterol\textquotedblright{}.) Say it once more.
\end{tcolorbox}
\section[\texorpdfstring{\ta{புரதம்} (puratam)}{புரதம் (puratam)}]{\ta{புரதம்} (puratam) — protein}

\label{lesson:TA-C487-puratam}



\begin{quote}
Before the new one: say the Tamil for to pull through, then the Tamil for body fat.
\end{quote}

\subsection*{You'll want to know: \ta{புரதம்}}
\textbf{\ta{புரதம்}} — \emph{puratam} — \textquotedblleft{}protein\textquotedblright{}.

\begin{quote}
\ta{பருப்பில்} \ta{நிறைய} \ta{புரதம்} \ta{இருக்கிறது}.
\end{quote}

\begin{quote}
\emph{paruppil niṟaiya puratam irukkiṟatu.}
\end{quote}

\begin{quote}
Lentils contain plenty of protein.
\end{quote}

A mass noun in science and nutrition.

\begin{cousinweb}
The origin of \textbf{\ta{புரதம்}} (\emph{puratam}), a modern scientific term, is not settled in our sources.
\end{cousinweb}

\begin{grammarlens}[title={What you've built}]
One more word for describing what you lived and felt.
\end{grammarlens}

\subsection*{Your turn}
\begin{itemize}
\raggedright
  \item \emph{Say it:} \emph{puratam}
  \item \emph{Say it:} \emph{puratam}, once more
  \item \emph{Say it:} \emph{paruppil niṟaiya puratam irukkiṟatu.} — Lentils contain plenty of protein.
  \item \emph{Recall:} say the Tamil for to pull through, then the Tamil for body fat, then say \emph{puratam} again
\end{itemize}

\begin{tcolorbox}[breakable,colback=teal!4,colframe=teal!35!black,title={Before you move on}]
What does \ta{புரதம்} mean? (\textquotedblleft{}Protein\textquotedblright{}.) Say it once more.
\end{tcolorbox}
\section[\texorpdfstring{\ta{வைட்டமின்} (vaiṭṭamiṉ)}{வைட்டமின் (vaiṭṭamiṉ)}]{\ta{வைட்டமின்} (vaiṭṭamiṉ) — a vitamin}

\label{lesson:TA-C487-vaittamin}



\begin{quote}
Before the new one: say the Tamil for body fat, then the Tamil for protein.
\end{quote}

\subsection*{You'll want to know: \ta{வைட்டமின்}}
\textbf{\ta{வைட்டமின்}} — \emph{vaiṭṭamiṉ} — \textquotedblleft{}a vitamin\textquotedblright{}.

\begin{quote}
\ta{பழங்களில்} \ta{நிறைய} \ta{வைட்டமின்கள்} \ta{உள்ளன}.
\end{quote}

\begin{quote}
\emph{paḻaṅkaḷil niṟaiya vaiṭṭamiṉkaḷ uḷḷaṉa.}
\end{quote}

\begin{quote}
Fruits contain plenty of vitamins.
\end{quote}

Plural: \ta{வைட்டமின்கள்}, \emph{vaiṭṭamiṉkaḷ}.

A loan ending in -\ta{ன்}, with the plural -\ta{கள்}.

\begin{cousinweb}
\textbf{\ta{வைட்டமின்}} (\emph{vaiṭṭamiṉ}) is a loan from English \emph{vitamin}.
\end{cousinweb}

\begin{grammarlens}[title={What you've built}]
One more word for describing what you lived and felt.
\end{grammarlens}

\subsection*{Your turn}
\begin{itemize}
\raggedright
  \item \emph{Say it:} \emph{vaiṭṭamiṉ}
  \item \emph{Say it:} \emph{vaiṭṭamiṉ}, once more
  \item \emph{Say it:} \emph{paḻaṅkaḷil niṟaiya vaiṭṭamiṉkaḷ uḷḷaṉa.} — Fruits contain plenty of vitamins.
  \item \emph{Recall:} say the Tamil for body fat, then the Tamil for protein, then say \emph{vaiṭṭamiṉ} again
\end{itemize}

\begin{tcolorbox}[breakable,colback=teal!4,colframe=teal!35!black,title={Before you move on}]
What does \ta{வைட்டமின்} mean? (\textquotedblleft{}A vitamin\textquotedblright{}.) Say it once more.
\end{tcolorbox}
\section[\texorpdfstring{\ta{செரி} (ceri)}{செரி (ceri)}]{\ta{செரி} (ceri) — to digest}

\label{lesson:TA-C487-ceri}



\begin{quote}
Before the new one: say the Tamil for protein, then the Tamil for a vitamin.
\end{quote}

\subsection*{You'll want to know: \ta{செரி}}
\textbf{\ta{செரி}} — \emph{ceri} — \textquotedblleft{}to digest\textquotedblright{}.

\begin{quote}
\ta{இரவு} \ta{உணவு} \ta{இன்னும்} \ta{செரிக்கவில்லை}.
\end{quote}

\begin{quote}
\emph{iravu uṇavu iṉṉum cerikkavillai.}
\end{quote}

\begin{quote}
Dinner has not been digested yet.
\end{quote}

Past (it …): \ta{செரித்தது}, \emph{cerittatu}.

The food is the subject, so the past is \ta{செரித்தது}. \ta{செரிமானம்}, digestion.

\begin{cousinweb}
\textbf{\ta{செரி}} (\emph{ceri}) is a Tamil verb whose origin is not settled; DEDR has no entry for it that we can point to.
\end{cousinweb}

\begin{grammarlens}[title={What you've built}]
One more word for describing what you lived and felt.
\end{grammarlens}

\subsection*{Your turn}
\begin{itemize}
\raggedright
  \item \emph{Say it:} \emph{ceri}
  \item \emph{Say it:} \emph{ceri}, once more
  \item \emph{Say it:} \emph{iravu uṇavu iṉṉum cerikkavillai.} — Dinner has not been digested yet.
  \item \emph{Recall:} say the Tamil for protein, then the Tamil for a vitamin, then say \emph{ceri} again
\end{itemize}

\begin{tcolorbox}[breakable,colback=teal!4,colframe=teal!35!black,title={Before you move on}]
What does \ta{செரி} mean? (\textquotedblleft{}To digest\textquotedblright{}.) Say it once more.
\end{tcolorbox}
\section[\texorpdfstring{\ta{ஏப்பம்} \ta{விடு} (ēppam viṭu)}{ஏப்பம் விடு (ēppam viṭu)}]{\ta{ஏப்பம்} \ta{விடு} (ēppam viṭu) — to burp; to belch}

\label{lesson:TA-C487-eppam-vitu}



\begin{quote}
Before the new one: say the Tamil for a vitamin, then the Tamil for to digest.
\end{quote}

\subsection*{You'll want to know: \ta{ஏப்பம்} \ta{விடு}}
\textbf{\ta{ஏப்பம்} \ta{விடு}} — \emph{ēppam viṭu} — \textquotedblleft{}to burp; to belch\textquotedblright{}.

\begin{quote}
\ta{நன்றாகச்} \ta{சாப்பிட்ட} \ta{பிறகு} \ta{தாத்தா} \ta{ஏப்பம்} \ta{விட்டார்}.
\end{quote}

\begin{quote}
\emph{naṉṟākac cāppiṭṭa piṟaku tāttā ēppam viṭṭār.}
\end{quote}

\begin{quote}
After a good meal grandfather burped.
\end{quote}

Past (I …): \ta{ஏப்பம்} \ta{விட்டேன்}, \emph{ēppam viṭṭēṉ}.

\ta{ஏப்பம்} \textquotedblleft{}a belch\textquotedblright{} + \ta{விடு} \textquotedblleft{}to let out\textquotedblright{}.

\begin{cousinweb}
\textbf{\ta{ஏப்பம்} \ta{விடு}} (\emph{ēppam viṭu}) is native: \ta{ஏப்பம்} \emph{ēppam}, \textquotedblleft{}a belch\textquotedblright{} (DEDR 897), with \ta{விடு} \textquotedblleft{}to let out\textquotedblright{}.
\end{cousinweb}

\begin{grammarlens}[title={What you've built}]
One more word for describing what you lived and felt.
\end{grammarlens}

\subsection*{Your turn}
\begin{itemize}
\raggedright
  \item \emph{Say it:} \emph{ēppam viṭu}
  \item \emph{Say it:} \emph{ēppam viṭu}, once more
  \item \emph{Say it:} \emph{naṉṟākac cāppiṭṭa piṟaku tāttā ēppam viṭṭār.} — After a good meal grandfather burped.
  \item \emph{Recall:} say the Tamil for a vitamin, then the Tamil for to digest, then say \emph{ēppam viṭu} again
\end{itemize}

\begin{tcolorbox}[breakable,colback=teal!4,colframe=teal!35!black,title={Before you move on}]
What does \ta{ஏப்பம்} \ta{விடு} mean? (\textquotedblleft{}To burp; to belch\textquotedblright{}.) Say it once more.
\end{tcolorbox}
